\documentclass[11pt,a4paper]{article}

\usepackage[margin=2.5cm]{geometry}
\usepackage{amsmath,amssymb}
\usepackage{booktabs}
\usepackage{graphicx}
\usepackage[colorlinks=true,linkcolor=blue,citecolor=blue,urlcolor=blue]{hyperref}

\newcommand{\DSigma}{\Delta\Sigma}
\newcommand{\Msun}{\mathrm{M}_\odot}
\newcommand{\dd}{\mathrm{d}}

\title{Destacking individual cluster lensing profiles versus the CoMaLit
programme:\\ a comparison of statistical strategies}
\author{A.~Z.~Vitorelli \and Claude (Anthropic)}
\date{July 15, 2026}

\begin{document}
\maketitle

\begin{abstract}
We compare two statistical programmes that address population-level galaxy
cluster relations from weak-lensing data: the CoMaLit series of Sereno,
Ettori and collaborators
\cite{comalit1,comalit2,comalit3,comalit4,comalit5,comalit6}, which
calibrates mass--observable scaling relations by Bayesian hierarchical
regression of compressed (point-estimate) masses and proxies; and the
\emph{destacking} approach developed in this repository, which fits a
concentration--observable relation jointly to the individual, un-stacked
excess surface density profiles of every cluster in the sample. The two
programmes occupy opposite ends of the data-compression axis: CoMaLit pairs
a maximal compression of the measurements with a fully stochastic population
model, while destacking pairs a minimal compression with a (currently)
deterministic population model. We show that the calibration failures
measured in destacking simulations at DES-like source density are the
profile-space manifestation of exactly the intrinsic-scatter effects that
CoMaLit-I identified in scaling-relation space, and we derive the
hierarchical extension of the destacking likelihood that imports the CoMaLit
population layer. We also identify what CoMaLit-style analyses cannot
recover at any price --- the concentration information residing in profile
shapes --- and summarise validation practices that each programme could
usefully adopt from the other.
\end{abstract}

\section{Introduction}

Galaxy cluster masses are not observables. Every use of clusters in
cosmology or structure-formation astrophysics therefore rests on relations
between mass (or halo structure) and measurable proxies --- richness,
Sunyaev--Zel'dovich (SZ) signal, X-ray luminosity, velocity dispersion,
weak-lensing shear --- together with a statistical model of how individual
clusters scatter about those relations and of how the sample was selected.

The \textbf{CoMaLit} (Comparing Masses in Literature) series
\cite{comalit1,comalit2,comalit3,comalit4,comalit5,comalit6} constitutes a
systematic programme for this problem, built on Bayesian hierarchical linear
regression in logarithmic variables, with latent true masses, intrinsic
scatter on all axes, redshift evolution of both relation and scatter, and
explicit treatment of Malmquist/Eddington selection effects. Its
methodological engine is distributed as the LIRA package \cite{lira}.

The \textbf{destacking} programme (this repository) approaches a sibling
problem --- the dependence of halo concentration on a formation-time proxy,
the magnitude gap $\Delta m$ between the brightest and second (or fourth)
brightest cluster galaxies --- by refusing the compression step altogether:
rather than reducing each cluster to a mass estimate, or stacking clusters
into high signal-to-noise bins, it forward-models the noisy $\DSigma$
profile of \emph{every individual cluster} and fits the relation parameters
jointly across the sample. A prototype comparison found the joint fit of
$500$ individual profiles delivered roughly twice tighter constraints on
$(\log M, c)$ than stacking the same clusters.

Sections~\ref{sec:comalit} and \ref{sec:destacking} state the two models in
their own notation. Section~\ref{sec:comparison} compares the strategies
dimension by dimension. Section~\ref{sec:lessons} derives what each can
import from the other, including the hierarchical extension of the
destacking likelihood suggested by the CoMaLit framework.

\section{The CoMaLit statistical model}
\label{sec:comalit}

\subsection{Scope of the series}

The six papers considered here contribute, in order: the demonstration that
intrinsic scatter in mass \emph{estimates} biases scaling relations if
ignored, with measured intrinsic scatters of ${\sim}15$ per cent for
weak-lensing and ${\sim}25$ per cent for hydrostatic masses and up to
${\sim}40$ per cent divergence between analyses of the same clusters by
different groups (CoMaLit-I, \cite{comalit1}); the calibration of the
Planck $Y_{500}$--$M_{500}$ relation and the operational distinction between
\emph{symmetric} and \emph{conditional} scaling relations (CoMaLit-II,
\cite{comalit2}); the standardized LC$^2$ literature catalogues of
weak-lensing masses (CoMaLit-III, \cite{comalit3}); the general
time-evolving hierarchical method with self-calibrated completeness
(CoMaLit-IV, \cite{comalit4}); mass forecasting with masses of uncalibrated
clusters treated as missing data (CoMaLit-V, \cite{comalit5}); and the
measurement of intrinsic proxy scatter from \emph{stacked} weak lensing by
modelling the occupancy dependence of the stacked scatter (CoMaLit-VI,
\cite{comalit6}).

\subsection{The regression model}
\label{sec:lira}

In the LIRA formulation \cite{lira,comalit4}, each cluster $i$ carries a
latent covariate $Z_i$ (e.g.\ the true log-mass, suitably normalised). The
latent proxy value $X_{Z,i}$ and latent response $Y_{Z,i}$ are linear in
$Z_i$ with a redshift-evolution term,
\begin{align}
X_{Z,i} &= \alpha_{X|Z} + \beta_{X|Z} Z_i + \gamma_{X|Z}\, T(z_i),
\label{eq:lira-x}\\
Y_{Z,i} &= \alpha_{Y|Z} + \beta_{Y|Z} Z_i + \gamma_{Y|Z}\, T(z_i),
\label{eq:lira-y}
\end{align}
where $T(z)$ is a chosen time function, e.g.\
$T(z)=\log\!\left[(1+z)/(1+z_{\rm ref})\right]$ or the log ratio of Hubble
parameters or distances. True values scatter about the latent relation with
intrinsic dispersions that may themselves evolve,
\begin{equation}
X_i \sim \mathcal{N}\!\big(X_{Z,i},\, \sigma_{X|Z}^2(z_i)\big),
\qquad
Y_i \sim \mathcal{N}\!\big(Y_{Z,i},\, \sigma_{Y|Z}^2(z_i)\big),
\label{eq:lira-scatter}
\end{equation}
and the observations add heteroscedastic, possibly correlated measurement
errors,
\begin{equation}
\begin{pmatrix} x_i \\ y_i \end{pmatrix}
\sim
\mathcal{N}\!\left(
\begin{pmatrix} X_i \\ Y_i \end{pmatrix},
\begin{pmatrix}
\delta_{x,i}^2 & \rho_i\,\delta_{x,i}\delta_{y,i}\\
\rho_i\,\delta_{x,i}\delta_{y,i} & \delta_{y,i}^2
\end{pmatrix}
\right).
\label{eq:lira-meas}
\end{equation}
The decisive ingredient is the prior on the latent population,
\begin{equation}
Z_i \sim \sum_k \pi_k\,
\mathcal{N}\!\big(\mu_{Z,k}(z_i),\, \sigma_{Z,k}^2(z_i)\big),
\label{eq:lira-pop}
\end{equation}
a (mixture of) Gaussian(s) whose mean and width may drift with redshift,
which stands in for the mass function convolved with the selection.
Selection on the response is modelled as truncation: if only objects with
$y_i > y_{{\rm th},i}$ enter the sample, the likelihood of each observed
object is renormalised by the detection probability,
\begin{equation}
p(y_i \mid Y_i) \;=\;
\frac{\mathcal{N}\big(y_i;\, Y_i,\, \delta_{y,i}^2\big)\,
\Theta(y_i - y_{{\rm th},i})}
{\displaystyle \int_{y_{{\rm th},i}}^{\infty}
\mathcal{N}\big(y';\, Y_i,\, \delta_{y,i}^2\big)\,\dd y'},
\label{eq:lira-malmquist}
\end{equation}
with the threshold (completeness function) itself self-calibrated in
CoMaLit-IV. Posterior sampling is by Gibbs (JAGS) since all conditionals
are (truncated) Gaussians.

Two consequences of this structure organise the whole series:

\paragraph{Slope attenuation.} If intrinsic scatter and measurement errors
on the covariate are ignored, the recovered slope is biased shallow by the
classical attenuation factor,
\begin{equation}
\big\langle \hat\beta \big\rangle \;\simeq\;
\beta \,
\frac{\sigma_Z^2}{\sigma_Z^2 + \sigma_{X|Z}^2 + \delta_x^2},
\label{eq:attenuation}
\end{equation}
and the apparent intrinsic scatter of the relation is correspondingly
inflated --- the central quantitative point of CoMaLit-I.

\paragraph{Symmetric versus conditional relations.} With scatter in the
model, the ``symmetric'' relation (the underlying
Eqs.~\ref{eq:lira-x}--\ref{eq:lira-y}) differs from the ``conditional''
expectation $E[Y \mid x]$ appropriate for forecasting or number counts; for
jointly Gaussian variables the conditional slope carries the same
attenuation form as Eq.~(\ref{eq:attenuation}). CoMaLit-II measures both
for the Planck clusters (finding, e.g., symmetric slopes $1.4$--$1.9$ but
conditional slopes $1.1$--$1.6$ for $Y_{500}$--$M_{500}$), and CoMaLit-V
builds mass forecasting on the conditional relation with unmeasured masses
treated as missing data estimated jointly with the regression parameters.

\paragraph{Stacked scatter rescaling.} CoMaLit-VI extends the framework to
stacked lensing: for a bin containing $N_k$ clusters the intrinsic
contribution to the scatter of the stacked signal is rescaled,
\begin{equation}
\sigma_{Y,{\rm stack},k}^2 \;\simeq\; \frac{\sigma_{Y|Z}^2}{N_k}
\;+\; \delta_{y,{\rm stack},k}^2,
\label{eq:comalit6}
\end{equation}
so that the \emph{individual} intrinsic scatter $\sigma_{Y|Z}$ remains
identifiable from stacked measurements provided bins of different occupancy
are available. Applied to AMICO clusters in KiDS-DR3, this yields richness
as a low-scatter mass proxy and $\sim 20$ per cent mass forecasts down to
$\sim 10^{13}\,\Msun$.

\section{The destacking model}
\label{sec:destacking}

\subsection{Forward model}

Destacking targets the concentration--magnitude-gap relation at fixed mass
proxy. For cluster $i$ with richness $\lambda_i$, redshift $z_i$ and
magnitude gap $\Delta m_i$, the model concentration is
\begin{equation}
\log_{10} c_i \;=\; \log_{10} c_0
\;+\; \alpha\,\big(\log_{10} M_i - \log_{10} M_{\rm piv}\big)
\;+\; \beta\,\log_{10}\!\frac{1+z_i}{1+z_{\rm ref}}
\;+\; \gamma\,\log_{10}\!\frac{\Delta m_i}{\Delta m_{\rm ref}},
\label{eq:crel}
\end{equation}
with the mass tied deterministically to richness through the McClintock
et al.\ relation \cite{mcclintock19} plus a single global calibration
offset,
\begin{equation}
\log_{10} M_i \;=\; \mu_{\rm MC}(\lambda_i, z_i) \;+\; \log_{10} M_{\rm norm}.
\label{eq:mlambda}
\end{equation}
The predicted profile is the centred/miscentred one-halo NFW
\cite{nfw97,wright00} excess surface density evaluated at the measured
radii $R_b$,
\begin{equation}
m_{i,b}(\theta) \;=\;
(1 - f_{{\rm mis},i})\,\DSigma_{\rm NFW}\big(R_b;\, M_i, c_i, z_i\big)
\;+\; f_{{\rm mis},i}\,\DSigma_{\rm mis}\big(R_b;\, M_i, c_i, z_i,
\sigma_{{\rm off},i}\big),
\label{eq:model}
\end{equation}
with $\theta = (c_0, \alpha, \beta, \gamma, \log_{10} M_{\rm norm})$ and the
two-halo term subtracted; a global amplitude correction $A$ is available in
the code but pinned to $A=1$ (its effect was found negligible in $\gamma$
scans).

\subsection{Per-cluster measurement}

Each cluster's data vector is the response-corrected metacalibration
\cite{sheldon17} estimator of $\DSigma$ in radial bins,
\begin{equation}
\widehat{\DSigma}_{i,b} \;=\;
\frac{\displaystyle\sum_{j \in b} w_j\, e_{t,j}\,\Sigma_{{\rm cr},j}^{-1}}
{\displaystyle\sum_{j \in b} w_j\, R_{t,j}\, \Sigma_{{\rm cr},j}^{-2}},
\label{eq:estimator}
\end{equation}
with $e_{t,j}$ the tangential ellipticity, $R_{t,j}$ the tangential
response, $w_j$ the weight and $\Sigma_{{\rm cr},j}$ the critical surface
density of source $j$. Both the data vector and its covariance
$\mathsf{C}_i$ are estimated from the \emph{same} galaxy-level bootstrap
(no boost correction, consistent with the boost-free model of
Eq.~\ref{eq:model}); bins in which the bootstrap fails to return a finite,
positive variance (too few sources) are masked per cluster.

\subsection{Joint likelihood and combination}

The relation parameters are constrained by the product of per-cluster
Gaussians over the valid bins $V_i$,
\begin{equation}
\ln \mathcal{L}(\theta) \;=\;
-\frac{1}{2}\sum_{i=1}^{N}
\Big[
\mathbf{r}_i^{\sf T} \mathsf{C}_i^{-1} \mathbf{r}_i
+ \ln\det \mathsf{C}_i
+ |V_i| \ln 2\pi
\Big],
\qquad
\mathbf{r}_i = \big(\widehat{\DSigma}_{i,b} - m_{i,b}(\theta)\big)_{b\in V_i},
\label{eq:jointlike}
\end{equation}
evaluated by Cholesky decomposition and sampled with the nautilus nested
sampler \cite{nautilus} under Gaussian priors on $\theta$. Sampling
robustness is addressed by cluster bootstrap: the sample of $N$ clusters is
resampled with replacement $B$ times, the fit repeated per resample, and
the quoted measurement is the equal-weight \emph{mixture} of the converged
bootstrap posteriors,
\begin{equation}
\tilde p(\gamma \mid D) \;=\;
\frac{1}{B_{\rm conv}} \sum_{b\,\in\,{\rm converged}} p_b(\gamma \mid D_b),
\qquad
\hat\gamma = \mathrm{median}\big[\tilde p\big],\quad
\sigma_\gamma = \tfrac{1}{2}\big(\tilde p_{84} - \tilde p_{16}\big).
\label{eq:mixture}
\end{equation}
The mixture width automatically combines the within-fit posterior width
with the between-resample scatter,
$\sigma_\gamma^2 \approx \sigma_{\rm within}^2 + \sigma_{\rm between}^2$,
and therefore remains honest even where the individual posteriors are
overconfident.

\subsection{Validation results (July 2026)}
\label{sec:validation}

The pipeline is validated end-to-end on synthetic surveys: triaxial NFW
shear injection with per-galaxy metacalibration responses, DES-Y3-like
source redshift distributions, Poisson-sampled richness with a
$\lambda > \lambda_{\rm th}$ selection, optional miscentering. A
$2\times2$ test matrix ($\{$DES, LSST source density$\}\times\{$triaxial +
miscentred, spherical + centred$\}$; $1500$ clusters; $50$ bootstrap fits
each; $\gamma_{\rm true}=+0.2$) gave the results of
Table~\ref{tab:matrix}. The pull is defined per bootstrap fit as
$(\hat\gamma_b - \gamma_{\rm true})/\sigma_b$ with $\sigma_b$ the per-fit
posterior half-width.

\begin{table}
\centering
\caption{Destacking validation matrix (July 2026). $\sigma_e = 0.27$
per component in all cells; densities in galaxies\,arcmin$^{-2}$.}
\label{tab:matrix}
\begin{tabular}{lcccc}
\toprule
 & DES (7.5), triax. & DES (7.5), spher. & LSST (40), triax. & LSST (40), spher.\\
\midrule
mean $\hat\gamma$ ($\gamma_{\rm true}=0.2$) & $+0.383\pm0.074$ & $+0.323\pm0.048$ & $+0.179\pm0.013$ & $+0.222\pm0.010$\\
scatter / median $\sigma_b$ & 2.58 & 2.70 & 1.09 & 1.22\\
pull RMS & 3.08 & 2.70 & 1.12 & 1.22\\
\bottomrule
\end{tabular}
\end{table}

At LSST density the method is essentially calibrated and unbiased, with
$\sigma_\gamma \simeq 0.06$--$0.08$ from only $1500$ clusters. At DES
density the per-fit posteriors are overconfident by a factor
${\sim}2.6$--$2.7$ and the point estimates over-dispersed, with a
${\sim}+0.14\pm0.04$ bias when the two DES cells are combined.

\section{Strategy comparison}
\label{sec:comparison}

\subsection{The compression axis}

Table~\ref{tab:compare} summarises the structural differences. The deepest
one is where each programme places the stochasticity. CoMaLit treats the
\emph{measurement} of each object as trivially summarised (a Gaussian with
known width) and the \emph{population} as richly stochastic
(Eqs.~\ref{eq:lira-scatter}--\ref{eq:lira-pop}). Destacking treats the
measurement as rich --- the full profile with its bootstrap covariance,
Eq.~(\ref{eq:jointlike}) --- and the population as deterministic: given
$(\lambda_i, z_i, \Delta m_i)$ and $\theta$, the model profile of
Eq.~(\ref{eq:model}) is a point, with no intrinsic scatter anywhere.

\begin{table}
\centering
\caption{Structural comparison.}
\label{tab:compare}
\begin{tabular}{p{4.0cm}p{5.1cm}p{5.1cm}}
\toprule
 & CoMaLit & destacking\\
\midrule
input datum per cluster & point estimate + error & $\DSigma$ vector + covariance\\
population model & latent $Z$, intrinsic scatter both axes, evolving (Eqs.~\ref{eq:lira-x}--\ref{eq:lira-pop}) & deterministic (Eqs.~\ref{eq:crel}--\ref{eq:model})\\
selection effects & truncation + population prior, self-calibrated completeness (Eq.~\ref{eq:lira-malmquist}) & present in simulations, absent from likelihood\\
target parameters & $\alpha, \beta$, evolution, $\sigma_{Y|Z}$ of mass--proxy relations & $\gamma$ (and $c_0,\alpha,\beta$) of the $c$--$\Delta m$ relation\\
inference & Gibbs sampling (JAGS/LIRA) & nested sampling (nautilus), cluster bootstrap, posterior mixture (Eq.~\ref{eq:mixture})\\
concentration information & integrated out at compression & the carrier of the signal\\
validation & literature cross-comparison, regression consistency & end-to-end synthetic surveys with shared code paths\\
\bottomrule
\end{tabular}
\end{table}

\subsection{The same disease in two spaces}

The DES-density pathologies of Section~\ref{sec:validation} are the
profile-space image of CoMaLit-I's central result. The destacking
likelihood, Eq.~(\ref{eq:jointlike}), asserts that each cluster's residual
is purely shape noise. In truth the residual also contains (i) the error of
the deterministic mass map, Eq.~(\ref{eq:mlambda}), driven by the Poisson
scatter of $\lambda$ at fixed $M$
($\sigma_{\ln M} \simeq |\partial \mu_{\rm MC}/\partial\ln\lambda| /
\sqrt{\lambda}$, i.e.\ $\sim 0.08$--$0.11$ dex at
$\lambda \sim 30$--$60$); (ii) any intrinsic scatter of $c$ about
Eq.~(\ref{eq:crel}); and (iii) profile scatter from orientation and
centring not captured by the mean model. Omitting these terms makes each
per-cluster likelihood too sharp; the joint posterior of $N$ such terms is
too sharp by roughly the same factor, and the estimator is correspondingly
over-dispersed --- pull RMS $\approx 3$ rather than $1$, exactly as
measured. The attenuation mechanism of Eq.~(\ref{eq:attenuation}) also has
a destacking analogue: scatter in the effective covariate
($\log \Delta m$ leverage against noisy per-cluster masses) biases the
recovered slope, and the sign and size of the measured DES-density bias
($+0.14\pm0.04$) should be reproducible from this mechanism --- a
quantitative check that has not yet been carried out.

\subsection{Stacking, destacking, and CoMaLit-VI}

CoMaLit-VI and destacking answer the same objection --- stacking erases
population variance --- in opposite ways. CoMaLit-VI keeps stacks and
recovers the individual intrinsic scatter through the occupancy scaling of
Eq.~(\ref{eq:comalit6}); it stays in the linear-Gaussian world, at the cost
of retaining binning choices and averaging away profile-shape information
within each bin. Destacking keeps clusters and pays two prices measured in
this repository: computational (each likelihood evaluation costs
$\mathcal{O}(N)$ Cholesky solves; wall-time truncation of the sampler at
$15000$ clusters produced collapsed posteriors until the budget was made
configurable) and statistical (the joint likelihood is fragile exactly
where per-cluster information is poor). The regime boundary is empirical:
at LSST density destacking is calibrated and reaches
$\sigma_\gamma\!\sim\!0.08$ with $1500$ clusters; at DES density its
uncorrected form is neither calibrated nor unbiased.

\subsection{What compression destroys}

The clearest asymmetry: $\gamma$ in Eq.~(\ref{eq:crel}) is not accessible
to CoMaLit-style regression of literature masses, because the
concentration of an individual cluster does not survive compression into
$\hat M$. The signal destacking targets lives in the systematic
$\Delta m$-dependence of profile \emph{shapes} at fixed mass proxy. In the
DES-density fits this is visible in the posteriors: $c_0, \alpha, \beta$
return their priors while $\gamma$ (broad prior) carries all the data
information --- the joint profile fit is, in effect, a measurement of
$\gamma$ alone. Conversely, everything CoMaLit needs is already in its
compressed inputs, and its linear-Gaussian structure buys exact
conditionals, transparent selection corrections, and trivial computing ---
none of which destacking currently enjoys.

\section{Lessons in both directions}
\label{sec:lessons}

\subsection{Importing the CoMaLit population layer into destacking}

The hierarchical completion of Eq.~(\ref{eq:jointlike}) follows the
CoMaLit template. Promote each cluster's mass to a latent variable with a
richness likelihood and a selection-truncated population prior, and give
the concentration relation an intrinsic scatter:
\begin{equation}
\mathcal{L}(\theta, \sigma_{\ln c}) \;=\;
\prod_{i=1}^{N} \int \dd M_i \int \dd c_i \;
p\big(\widehat{\mathbf{D}}_i \,\big|\, m_i(M_i, c_i)\big)\;
\mathcal{N}\big(\ln c_i;\, \ln c(M_i; \theta),\, \sigma_{\ln c}^2\big)\;
\frac{P(\lambda_i \mid M_i)\; n(M_i, z_i)}
{P(\lambda > \lambda_{\rm th} \mid z_i)},
\label{eq:hierarchical}
\end{equation}
with $P(\lambda \mid M)$ Poisson or log-normal and $n(M, z)$ a mass
function (or, following Eq.~\ref{eq:lira-pop}, an effective Gaussian
mixture). A cheaper first step that captures most of the variance
linearises the latent integrals into an additive model-error term in each
covariance,
\begin{equation}
\mathsf{C}_i \;\rightarrow\;
\mathsf{C}_i
+ \sigma_{\ln c}^2\, \mathbf{g}_i \mathbf{g}_i^{\sf T}
+ \sigma_{\ln M|\lambda}^2\, \mathbf{h}_i \mathbf{h}_i^{\sf T},
\qquad
\mathbf{g}_i = \frac{\partial \mathbf{m}_i}{\partial \ln c},\quad
\mathbf{h}_i = \frac{\partial \mathbf{m}_i}{\partial \ln M},
\label{eq:modelerror}
\end{equation}
which keeps Eq.~(\ref{eq:jointlike}) Gaussian and nearly free. Ranked by
expected impact on the measured DES-density problems:
\begin{enumerate}
\item intrinsic and proxy scatter, via Eq.~(\ref{eq:modelerror}) first and
Eq.~(\ref{eq:hierarchical}) eventually --- the direct cure for the
$2.6\times$ overconfidence, and plausibly for part of the bias; broader
per-cluster likelihoods should also soften the sampler-hostile geometry;
\item selection truncation, the denominator of Eq.~(\ref{eq:hierarchical})
--- the Eddington bias near $\lambda_{\rm th}$ is currently in the sample
but not in the model;
\item the symmetric/conditional distinction of CoMaLit-II when quoting
$\gamma$ for a selected sample;
\item CoMaLit-V missing-data treatment for clusters without measured
$\Delta m$ (the current convention --- no $\gamma$ correction, i.e.\
$\log(\Delta m/\Delta m_{\rm ref}) \equiv 0$ --- is a hidden informative
choice);
\item time-evolving scatter once items 1--2 exist.
\end{enumerate}

\subsection{Importing destacking practices into CoMaLit-style analyses}

\begin{enumerate}
\item \emph{Profile-level joint inference where depth allows.} At
LSST-like density the compression to point masses is no longer forced;
concentration and mass can be constrained jointly per cluster within one
homogeneous model, removing the between-group heterogeneity that
CoMaLit-I quantifies at up to $40$ per cent.
\item \emph{End-to-end synthetic-survey validation with shared code
paths.} The class of error exemplified by the metacalibration response
mismatch found in this repository --- the simulation injected unit-response
shear while the estimator of Eq.~(\ref{eq:estimator}) divided by
$R_t \simeq 0.82$, inflating every profile by ${\sim}18$--$22$ per cent ---
is invisible to regression-level consistency checks by construction, and
was caught within hours by a zero-shape-noise single-cluster test against
the forward model.
\item \emph{The bootstrap posterior mixture}, Eq.~(\ref{eq:mixture}), as a
cheap honesty device whenever there is reason to suspect a likelihood is
too sharp: the mixture width carries the between-resample variance without
any model extension.
\end{enumerate}

\section{Summary}

CoMaLit and destacking are complementary solutions to the same class of
problem, differing in where they spend their model complexity. CoMaLit's
hierarchical population layer --- intrinsic scatter, latent variables,
selection truncation --- is precisely what destacking's likelihood lacks,
and the calibration failures measured in destacking simulations at
DES-like density are the textbook symptoms (CoMaLit-I) of that lack.
Destacking's profile-level forward model is precisely what CoMaLit's
compressed inputs cannot express, and it is the only route to the
concentration--magnitude-gap slope $\gamma$ that motivates the programme.
Equation~(\ref{eq:hierarchical}) --- the destacking measurement model under
the CoMaLit population model --- is strictly more expressive than either
programme alone; its computational cost, against the already-demonstrated
fragility of the plain joint likelihood at low per-cluster signal-to-noise,
is the main open engineering question, and is most naturally answered in
the LSST-like regime where the joint likelihood is demonstrably
well-behaved.

\begin{thebibliography}{99}

\bibitem{comalit1}
Sereno M., Ettori S., 2015, MNRAS, 450, 3633
(CoMaLit-I: Bias and scatter in weak lensing and X-ray mass estimates of
clusters), \href{https://arxiv.org/abs/1407.7868}{arXiv:1407.7868}

\bibitem{comalit2}
Sereno M., Ettori S., 2015, MNRAS, 450, 3649
(CoMaLit-II: The scaling relation between mass and Sunyaev--Zel'dovich
signal for Planck selected galaxy clusters),
\href{https://arxiv.org/abs/1407.7869}{arXiv:1407.7869}

\bibitem{comalit3}
Sereno M., 2015, MNRAS, 450, 3665
(CoMaLit-III: Literature catalogues of weak lensing clusters of galaxies,
LC$^2$), \href{https://arxiv.org/abs/1409.5435}{arXiv:1409.5435}

\bibitem{comalit4}
Sereno M., Ettori S., 2015, MNRAS, 450, 3675
(CoMaLit-IV: Evolution and self-similarity of scaling relations with the
galaxy cluster mass),
\href{https://arxiv.org/abs/1502.05413}{arXiv:1502.05413}

\bibitem{comalit5}
Sereno M., Ettori S., 2017, MNRAS
(CoMaLit-V: Mass forecasting with proxies --- method and application to
weak lensing calibrated samples),
\href{https://arxiv.org/abs/1603.06581}{arXiv:1603.06581}

\bibitem{comalit6}
Sereno M., et al., 2020, MNRAS
(CoMaLit-VI: Intrinsic scatter in stacked relations --- the weak lensing
AMICO galaxy clusters in KiDS-DR3),
\href{https://arxiv.org/abs/2006.16264}{arXiv:2006.16264}

\bibitem{lira}
Sereno M., 2016, MNRAS, 455, 2149
(A Bayesian approach to linear regression in astronomy: the LIRA package),
\href{https://arxiv.org/abs/1509.05778}{arXiv:1509.05778}

\bibitem{nfw97}
Navarro J.~F., Frenk C.~S., White S.~D.~M., 1997, ApJ, 490, 493

\bibitem{wright00}
Wright C.~O., Brainerd T.~G., 2000, ApJ, 534, 34

\bibitem{mcclintock19}
McClintock T., et al.\ (DES Collaboration), 2019, MNRAS, 482, 1352
(Dark Energy Survey Year 1 results: weak lensing mass calibration of
redMaPPer galaxy clusters),
\href{https://arxiv.org/abs/1805.00039}{arXiv:1805.00039}

\bibitem{sheldon17}
Sheldon E.~S., Huff E.~M., 2017, ApJ, 841, 24
(Practical weak-lensing shear measurement with metacalibration)

\bibitem{nautilus}
Lange J.~U., 2023, MNRAS, 525, 3181
(nautilus: boosting Bayesian importance nested sampling with deep
learning), \href{https://arxiv.org/abs/2306.16923}{arXiv:2306.16923}

\bibitem{dutton14}
Dutton A.~A., Macci\`o A.~V., 2014, MNRAS, 441, 3359
(Cold dark matter haloes in the Planck era: evolution of structural
parameters for Einasto and NFW profiles)

\end{thebibliography}

\end{document}
